\chapter{Bicircle curves and their error vectors}
\label{chap:bicircle}

% archon:covers Gluck/Bicircle.lean

This chapter analyses the geometry of the curves reconstructed from a four-arc
step curvature $\kappa_0$ (\cref{def:four_arc_step_function}). Such a curve is a
\emph{bicircle}: four circular arcs, two cut from a circle of curvature $a$ and
two from a circle of curvature $b$, joined so that the tangent line turns
continuously through $2\pi$. The main result is Dahlberg's explicit formula
(Proposition 9.1 of the source) for the error vector of a bicircle, which is the
input to the winding argument of \cref{chap:winding}.

% SOURCE: [DeTurck-Gluck], §9, p. 15-16, lines 276-298 (read from references/deturck-gluck-fourvertex.txt)
% SOURCE QUOTE: "PROPOSITION 9.1. A plane curve, built with four arcs cut in
% alternation from two different size circles, and assembled so that the tangent
% line turns continuously through an angle of 2π , will close up if and only if
% opposite arcs are equal in length. ... Proof. We parametrize such a curve by
% the angle of inclination θ of its tangent line, suppose that 0 < a < b are the
% curvatures of the two circles to be used, and record the curvature function
% κ(θ) in the figure below. ... If α(θ) is the resulting curve, then its error
% vector is given by  E = α(2π) — α(0) = ∫02π eiθ ds = ∫02π eiθ ds/dθ dθ = ∫02π
% eiθ / κ(θ) dθ . We evaluate this explicitly as a sum of four integrals, and get
% E = (1/ia — 1/ib) ( [exp(iθ2) — exp(iθ1)] + [exp(iθ4) — exp(iθ3)] ) ."
\begin{definition}

\leanok
[Bicircle error vector]
  \label{def:bicircle_error_vector}
  \lean{Gluck.bicircleErrorVector}
  \uses{def:four_arc_step_function, def:error_vector}
  \textit{Source: DeTurck--Gluck, Proposition 9.1 (§9).}
  Fix $0 < a < b$ and breakpoints $0 \le \theta_1 < \theta_2 < \theta_3 <
  \theta_4 < 2\pi$, and let $\kappa_0$ be the four-arc step curvature
  (\cref{def:four_arc_step_function}) taking the value $b$ on
  $[\theta_1,\theta_2)\cup[\theta_3,\theta_4)$ and $a$ elsewhere. The
  \emph{bicircle error vector} is $E_{\mathrm{bi}}(a,b,\theta_1,\dots,\theta_4) =
  E(1/\kappa_0) = \int_0^{2\pi} e^{i\theta}/\kappa_0(\theta)\,d\theta \in
  \mathbb{C}$.
\end{definition}

% SOURCE QUOTE PROOF: see the §9 quote above — "We evaluate this explicitly as a
% sum of four integrals, and get E = (1/ia — 1/ib) ( [exp(iθ2) — exp(iθ1)] +
% [exp(iθ4) — exp(iθ3)] )."
% SIGN CONVENTION: the source writes the scalar as $(1/(ia)-1/(ib))$ for its
% arc convention. With our `stepCurvature` convention (value $b$ on
% $[\theta_1,\theta_2)\cup[\theta_3,\theta_4)$, value $a$ on the complement) the
% correct scalar is the negative, $(1/(ib)-1/(ia))$ — see the proof's term
% collection below. We state the form our convention actually produces (the one
% the Lean proof verifies); the two forms differ only by which curvature labels
% the $b$-arcs, and the downstream winding argument needs only a nonzero scalar
% multiple of the chord sum, so the sign is immaterial there.
\begin{lemma}

\leanok
[Explicit error vector of a bicircle]
  \label{lem:error_vector_of_bicircle}
  \lean{Gluck.bicircleErrorVector_eq}
  \uses{def:bicircle_error_vector}
  \textit{Source: DeTurck--Gluck, Proposition 9.1 (§9), error-vector formula.}
  With the notation of \cref{def:bicircle_error_vector},
  \[
    E_{\mathrm{bi}}(a,b,\theta_1,\theta_2,\theta_3,\theta_4)
      \;=\; \Bigl(\tfrac{1}{ib} - \tfrac{1}{ia}\Bigr)
            \Bigl(\bigl[e^{i\theta_2} - e^{i\theta_1}\bigr]
                + \bigl[e^{i\theta_4} - e^{i\theta_3}\bigr]\Bigr).
  \]
\end{lemma}
\begin{proof}
  \uses{def:bicircle_error_vector, lem:bicircle_arc_integral, lem:bicircle_arc_integrable}
  \leanok
  The weight $\rho_0 = 1/\kappa_0$ is piecewise constant: it equals $1/b$ on the
  two arcs $[\theta_1,\theta_2)$ and $[\theta_3,\theta_4)$ and $1/a$ on the two
  complementary arcs. Split the integral $E(\rho_0) = \int_0^{2\pi}
  e^{i\theta}\rho_0(\theta)\,d\theta$ into the four arcs. On each arc $\rho_0$ is
  constant and $\int_{c}^{d} e^{i\theta}\,d\theta = \frac{1}{i}\bigl(e^{id} -
  e^{ic}\bigr)$. Summing,
  \[
    E(\rho_0)
      = \tfrac{1}{ib}\bigl(e^{i\theta_2}-e^{i\theta_1}\bigr)
      + \tfrac{1}{ia}\bigl(e^{i\theta_3}-e^{i\theta_2}\bigr)
      + \tfrac{1}{ib}\bigl(e^{i\theta_4}-e^{i\theta_3}\bigr)
      + \tfrac{1}{ia}\bigl(e^{i(\theta_1+2\pi)}-e^{i\theta_4}\bigr).
  \]
  Using $e^{i(\theta_1+2\pi)} = e^{i\theta_1}$ and collecting the $1/(ia)$ and
  $1/(ib)$ terms, the four $1/(ia)$ contributions sum to
  $\tfrac{1}{ia}\bigl((e^{i\theta_3}-e^{i\theta_2}) +
  (e^{i\theta_1}-e^{i\theta_4})\bigr) = -\tfrac{1}{ia}\bigl((e^{i\theta_2} -
  e^{i\theta_1}) + (e^{i\theta_4} - e^{i\theta_3})\bigr)$, while the two $1/(ib)$
  contributions sum to $\tfrac{1}{ib}\bigl((e^{i\theta_2} - e^{i\theta_1}) +
  (e^{i\theta_4} - e^{i\theta_3})\bigr)$. Hence the total is exactly
  $\bigl(\tfrac{1}{ib} - \tfrac{1}{ia}\bigr)\bigl([e^{i\theta_2} - e^{i\theta_1}]
  + [e^{i\theta_4} - e^{i\theta_3}]\bigr)$, the stated form. (The source's
  $\bigl(\tfrac{1}{ia} - \tfrac{1}{ib}\bigr)$ is the same up to the overall sign,
  which only reflects the opposite labelling of the two curvature values on the
  arcs.)
\end{proof}

% --- Computational helpers for the four-arc integral (Lean-side; small lemmas
% that the proof of lem:error_vector_of_bicircle factors through). ---
\begin{lemma}

\leanok
[Arc integral of $e^{i\theta}$]
  \label{lem:integral_cexp_I}
  \lean{Gluck.integral_cexp_I}
  Project-bespoke. For real $c \le d$, $\int_c^d e^{i\theta}\,d\theta =
  -i\bigl(e^{id} - e^{ic}\bigr)$.
\end{lemma}
\begin{proof}

\leanok
  Fundamental theorem of calculus with antiderivative $-i\,e^{i\theta}$ (whose
  derivative is $e^{i\theta}$, using $i\cdot i = -1$).
\end{proof}

\begin{lemma}

\leanok
[Step integrand is a.e.\ constant on an arc]
  \label{lem:bicircle_arc_integrand_eq}
  \lean{Gluck.bicircle_arc_integrand_eq}
  \uses{def:four_arc_step_function}
  Project-bespoke. On an arc $(c,d]$ on which $\kappa_0 \equiv v$, the integrand
  $e^{i\theta}\rho_0(\theta)$ agrees almost everywhere with $e^{i\theta}(1/v)$
  (the two single breakpoints are null).
\end{lemma}
\begin{proof}

\leanok
  $\rho_0 = 1/\kappa_0 = 1/v$ on the half-open arc; the endpoint where it may
  differ is a single point, of Lebesgue measure zero.
\end{proof}

\begin{lemma}

\leanok
[Per-arc integral value]
  \label{lem:bicircle_arc_integral}
  \lean{Gluck.bicircle_arc_integral}
  \uses{lem:integral_cexp_I, lem:bicircle_arc_integrand_eq}
  Project-bespoke. On an arc $[c,d]$ where $\kappa_0 \equiv v$,
  $\int_c^d e^{i\theta}\rho_0\,d\theta = (1/v)\,(-i)\,(e^{id}-e^{ic})$.
\end{lemma}
\begin{proof}

\leanok
  Replace the integrand by its a.e.\ representative
  (\cref{lem:bicircle_arc_integrand_eq}), pull out the constant $1/v$, and apply
  \cref{lem:integral_cexp_I}.
\end{proof}

\begin{lemma}

\leanok
[Per-arc integrability]
  \label{lem:bicircle_arc_integrable}
  \lean{Gluck.bicircle_arc_integrable}
  \uses{def:four_arc_step_function}
  Project-bespoke. The integrand $e^{i\theta}\rho_0$ is interval-integrable on
  each of the five constant arcs.
\end{lemma}
\begin{proof}

\leanok
  On each arc the integrand agrees a.e.\ with a continuous function (the constant
  weight times $e^{i\theta}$), which is interval-integrable; integrability
  transfers along a.e.\ equality.
\end{proof}

\begin{corollary}

\leanok
[Bicircle closes iff opposite arcs equal]
  \label{cor:bicircle_closes_iff}
  \lean{Gluck.bicircleErrorVector_eq_zero_iff}
  \uses{lem:error_vector_of_bicircle, lem:closes_iff}
  \textit{Source: DeTurck--Gluck, Proposition 9.1 (closure criterion).}
  Since $a \ne b$ the scalar $\tfrac{1}{ia} - \tfrac{1}{ib} \ne 0$, so the
  bicircle closes up ($E_{\mathrm{bi}} = 0$) if and only if
  $\bigl(e^{i\theta_2} - e^{i\theta_1}\bigr) + \bigl(e^{i\theta_4} -
  e^{i\theta_3}\bigr) = 0$, i.e. the two chords $e^{i\theta_2} - e^{i\theta_1}$
  and $e^{i\theta_4} - e^{i\theta_3}$ are equal and opposite. Geometrically this
  is the source's ``opposite arcs are equal in length'': the arc-length of each
  arc is its angular extent divided by its curvature, and the condition equates
  the two pairs of opposite arcs.
\end{corollary}
\begin{proof}
  \uses{lem:error_vector_of_bicircle}
  \leanok
  Immediate from \cref{lem:error_vector_of_bicircle}: a nonzero complex scalar
  times a vector is zero iff the vector is zero.
\end{proof}
